\chapter{Wprowadzenie}
\label{chap:wprowadzenie}
\begin{tikzpicture}[
    node distance=9mm,
    startstop/.style={ellipse, draw, align=center, minimum width=2.6cm},
    process/.style={rectangle, draw, align=center, minimum width=3.4cm},
    decision/.style={diamond, draw, align=center, aspect=2.4, inner sep=1pt},
    every path/.style={-{Stealth}}
]
    \node[startstop] (start) {Start};
    \node[process, below=of start] (pobierz) {Pobierz tablicę};
    \node[process, below=of pobierz] (przebieg) {Wykonaj pełny przebieg\\porównań i zmian};
    \node[decision, below=14mm of przebieg] (warunek) {Czy wykonano\\zamianę?};
    \node[startstop, below=14mm of warunek] (koniec) {Koniec};
    \draw (start) -- (pobierz);
    \draw (pobierz) -- (przebieg);
    \draw (przebieg) -- (warunek);
    \draw (warunek) -- node[right] {Nie} (koniec);
    \draw (warunek.east) -- ++(1.5,0) |- (przebieg.east)
        node[pos=0.5, right, yshift=-15mm] {Tak};
\end{tikzpicture}

\section{Wprowadzenie}

Sortowanie polega na uporządkowaniu elementów zbioru według przyjętego
\textbf{kryterium}, np. \textit{rosnąco} lub malejąco.
Algorytmy sortowania wykorzystuje się podczas przetwarzania i wyszukiwania danych.
Jedną z prostych metod jest \textbf{sortowanie bąbelkowe}.

\section{Sortowanie bąbelkowe}

    \subsection{Zasada działania}

    Sortowanie bąbelkowe polega na porównywaniu par sąsiednich elementów tablicy.
    Jeśli element o indeksie \(i\) jest większy niż kolejny element, ich miejsca
    są zamieniane. Warunek zamiany przedstawia zależność \(a_i > a_{i+1}\).

    Kolejne przebiegi algorytmu powodują stopniowe przemieszczanie największych
    elementów na koniec tablicy.

    \subsection{Etapy algorytmu}

    \begin{enumerate}
        \item Porównaj dwa sąsiednie elementy tablicy.
        \item Zamień je miejscami, jeśli ich kolejność jest niewłaściwa.
        \item Przejdź do kolejnej pary elementów.
        \item Powtarzaj przebiegi aż do uporządkowania danych.
    \end{enumerate}
    Dwie cechy algorytmu:
    \begin{itemize}
        \item prostota implementacji,
        \item niska efektywność dla dużych zbiorów danych.
    \end{itemize}

\section{Złożoność obliczeniowa}

Dla typowej wersji sortowania bąbelkowego liczba porównań wykonywanych
w najgorszym przypadku wynosi:
\begin{equation}
    C(n) = \frac{n(n-1)}{2}
\end{equation}
Złożoność czasowa w najgorszym przypadku to \(O(n^2)\).
a dodatkowa złożoność pamięciowa wynosi \(O(1)\).
\section{Podstawowe zapisy matematyczne}
Suma elementów tablicy wynosi:
\[
    S = \sum_{i=1}^{n} a_i
\]

Funkcja i jej pochodna:
\[
    f(x) = x^3 + 2x^2 -5x, \qquad f'(x) = 3x^2 + 4x - 5
\]
\section{Wnioski}

Sortowanie bąbelkowe jest proste w implementacji i pomaga zrozumieć działanie
algorytmów porównawczych. Ze względu na kwadratową złożoność czasową nie jest
jednak odpowiednim wyborem do sortowania bardzo dużych zbiorów danych.
W arkuszu \textbf{Excel} sortowanie układa wiersze według wartości wybranej kolumny.